\section{Conclusion}

This thesis developed and evaluated a symbolic-music generation pipeline based
on REMI+-style MIDI tokenization, autoregressive language modelling, and
generated-MIDI evaluation. The work began from an xLSTM-focused setup but was
extended into a controlled comparison between xLSTM, a GPT-2-style Transformer,
and a LLaMA-style Transformer. The final repository now covers dataset
construction, leakage-aware splitting, preprocessing, tokenizer selection,
model training, checkpoint selection, W\&B logging, generated-MIDI scoring, and
interactive Gradio applications for both recurrent and Transformer checkpoints.

The strongest next-token model in the matched comparison is the selected xLSTM
checkpoint. Under the shared tokenizer, context length, batch size, and update
budget, it achieves the best held-out loss, perplexity, and token accuracy. The
Transformer baselines do not match that token-level performance, but they remain
useful comparison points: the LLaMA-style model is the stronger Transformer
baseline, while the GPT-2-style model obtains the lowest FMD in the generation
evaluation. The generated-MIDI metrics therefore show that musical output
quality cannot be reduced to one scalar metric or to validation loss alone.

The practical conclusion is that the selected xLSTM is the best model in this
thesis when next-token modelling and rhythmic-density agreement are prioritized,
but it is also the most expensive model to sample from. The Transformer models
are weaker language models in this setup but are faster generators and preserve
some generated-MIDI statistics competitively. This creates a real deployment
tradeoff: the preferred model depends on whether the priority is predictive
accuracy, generation speed, or a particular musical descriptor.

The Gradio applications complete the project by making the trained checkpoints
usable outside the batch-evaluation pipeline. They allow MIDI prompts,
checkpoint selection, sampling controls, iterative continuation, playback, and
diagnostics for both model families. They do not replace formal listening
evaluation, but they provide a practical inspection interface for future
qualitative work.

Overall, the thesis contributes a reproducible comparison workflow for symbolic
music generation rather than only a single trained model. It shows how tokenizer
choice, checkpoint selection, architecture, generation settings, and evaluation
metrics interact in practice. The resulting evidence supports xLSTM as the
strongest model under the current controlled setup, while also identifying clear
future directions: larger-data training, broader Transformer tuning, stricter
generation-length controls, and a structured human listening study.
